% 4 邻域、8 邻域与 m 邻接的教学示意；宽约105mm。
\begin{tikzpicture}[x=5.5mm,y=5.5mm,font=\figurefont\small]
  \tikzset{cell/.style={draw=BookRule,line width=0.5pt},
    center/.style={fill=FigureRed,draw=FigureRed},
    nb/.style={fill=FigureBlueFill,draw=FigureBlue,line width=0.6pt},
    on/.style={fill=FigureBlue,draw=FigureBlue}}
  % 4 邻域
  \foreach \x in {0,1,2}{\foreach \y in {0,1,2}{\draw[cell] (\x,\y) rectangle ++(1,1);}}
  \foreach \x/\y in {1/2,0/1,2/1,1/0}{\filldraw[nb] (\x,\y) rectangle ++(1,1);}
  \filldraw[center] (1,1) rectangle ++(1,1);
  \node at (1.5,-0.7) {$N_4(p)$};
  % 8 邻域
  \begin{scope}[xshift=22mm]
    \foreach \x in {0,1,2}{\foreach \y in {0,1,2}{\filldraw[nb] (\x,\y) rectangle ++(1,1);}}
    \filldraw[center] (1,1) rectangle ++(1,1);
    \node at (1.5,-0.7) {$N_8(p)$};
  \end{scope}
  % 8 邻接造成的多重路径
  \begin{scope}[xshift=44mm]
    \foreach \x in {0,1,2}{\foreach \y in {0,1,2}{\draw[cell] (\x,\y) rectangle ++(1,1);}}
    \foreach \x/\y in {1/2,2/2,1/1,0/0}{\filldraw[on] (\x,\y) rectangle ++(1,1);}
    \draw[FigureRed,line width=1pt] (1.5,2.5) -- (2.5,2.5);
    \draw[FigureRed,line width=1pt] (1.5,1.5) -- (2.5,2.5);
    \draw[FigureRed,line width=1pt] (1.5,2.5) -- (1.5,1.5);
    \draw[FigureRed,line width=1pt] (1.5,1.5) -- (0.5,0.5);
    \node at (1.5,-0.7) {8 邻接};
  \end{scope}
  % m 邻接去掉对角冗余
  \begin{scope}[xshift=66mm]
    \foreach \x in {0,1,2}{\foreach \y in {0,1,2}{\draw[cell] (\x,\y) rectangle ++(1,1);}}
    \foreach \x/\y in {1/2,2/2,1/1,0/0}{\filldraw[on] (\x,\y) rectangle ++(1,1);}
    \draw[FigureRed,line width=1pt] (1.5,2.5) -- (2.5,2.5);
    \draw[FigureRed,line width=1pt] (1.5,2.5) -- (1.5,1.5);
    \draw[FigureRed,line width=1pt] (1.5,1.5) -- (0.5,0.5);
    \node at (1.5,-0.7) {m 邻接};
  \end{scope}
\end{tikzpicture}
